\documentclass[11pt,a4paper]{article}
\usepackage[italian]{babel}
\usepackage[margin=2.5cm]{geometry}
\usepackage{parskip}
\usepackage{amsmath, amssymb, amsthm}
\usepackage{booktabs}
\usepackage{array}
\usepackage{xcolor}
\usepackage{needspace}
\usepackage{enumitem}
\usepackage{listings}
\usepackage{tikz}
\usepackage[most]{tcolorbox}
\usepackage[hidelinks]{hyperref}
\usetikzlibrary{decorations.pathreplacing, shapes.geometric, arrows.meta, positioning}

% Niente vedove/orfane: i paragrafi non lasciano righe isolate a fine/inizio pagina
\widowpenalty=10000
\clubpenalty=10000

% Elenchi più compatti
\setlist{itemsep=2pt, parsep=0pt, topsep=4pt, partopsep=0pt}

% --- Riquadri con bordo nero, senza numero (si spezzano tra due pagine) ---
% Uso: \begin{definizione} ... \end{definizione}
%      \begin{teorema}[Nome] ... \end{teorema}  (il nome è facoltativo)
\tcbset{nota/.style={enhanced jigsaw, breakable, colback=white,
  colframe=black, boxrule=0.5pt, sharp corners}}
\NewTColorBox{definizione}{}{nota,
  before upper={\textbf{Definizione.}\ }}
\NewTColorBox{teorema}{o}{nota,
  before upper={\textbf{Teorema\IfValueT{#1}{ (#1)}.}\ }}
\NewTColorBox{proposizione}{o}{nota,
  before upper={\textbf{Proposizione\IfValueT{#1}{ (#1)}.}\ }}
\NewTColorBox{proprieta}{o}{nota,
  before upper={\textbf{Proprietà\IfValueT{#1}{ (#1)}.}\ }}
\NewTColorBox{assioma}{o}{nota, boxrule=1pt,
  before upper={\textbf{Assioma\IfValueT{#1}{ (#1)}.}\ }}

% --- Ambienti semplici (senza riquadro) ---
% Il titolo sta in cima al blocco, anche se il contenuto inizia con un riquadro o
% con codice; e non resta mai da solo in fondo a una pagina.
% Uso: \begin{esempio}[Titolo facoltativo] ... \end{esempio}
\newcommand{\newrunin}[2]{%
  \NewDocumentEnvironment{#1}{o}{%
    \par\Needspace{4\baselineskip}\addvspace{\medskipamount}\noindent
    \textbf{#2}\IfValueT{##1}{ (##1)}\textbf{.}\ \ignorespaces}%
  {\par\addvspace{\medskipamount}}}
\newrunin{esempio}{Esempio}
\newrunin{esercizio}{Esercizio}
\newrunin{osservazione}{Osservazione}

% V e F nelle tavole di verità
\newcommand{\V}{\textsf{V}}
\newcommand{\F}{\textsf{F}}

% --- Codice C ---
% Uso: \begin{codice} ... \end{codice}      codice C numerato
%      \begin{comando} ... \end{comando}    comandi da terminale
%      \begin{risultato} ... \end{risultato} output di un programma
%      \lstinline|printf("ciao");|           codice dentro al testo
\lstdefinestyle{c}{
  language=C,
  basicstyle=\ttfamily\small,
  keywordstyle=\color{blue}\bfseries,
  commentstyle=\color{gray}\itshape,
  stringstyle=\color{red!70!black},
  numbers=left,
  numberstyle=\tiny\color{gray},
  numbersep=8pt,
  columns=flexible,
  keepspaces=true,
  breaklines=true,
  showstringspaces=false,
  tabsize=4,
  morekeywords={include, define},
  % lettere accentate nei commenti
  literate={à}{{\`a}}1 {è}{{\`e}}1 {é}{{\'e}}1 {ì}{{\`i}}1
           {ò}{{\`o}}1 {ù}{{\`u}}1 {È}{{\`E}}1 {À}{{\`A}}1
}
\lstset{style=c}
\newtcblisting{codice}{listing only, nota, left=6mm,
  listing options={style=c}}
\newtcblisting{comando}{listing only, nota,
  listing options={basicstyle=\ttfamily\small, numbers=none, language={},
    columns=flexible, keepspaces=true, breaklines=true}}
\newtcblisting{risultato}{listing only, enhanced jigsaw, breakable,
  colback=black!5, colframe=black, boxrule=0.5pt, sharp corners,
  listing options={basicstyle=\ttfamily\small, numbers=none, language={},
    columns=flexible, keepspaces=true, breaklines=true}}

% --- Diagrammi di flusso (simboli standard) ---
\tikzset{
  terminale/.style = {ellipse, draw, minimum width=2.5cm, minimum height=0.9cm},
  io/.style        = {trapezium, trapezium left angle=70, trapezium right angle=110,
                      draw, minimum width=2.5cm, minimum height=0.9cm},
  processo/.style  = {rectangle, draw, minimum width=2.5cm, minimum height=0.9cm},
  decisione/.style = {diamond, draw, aspect=2, minimum width=2.5cm},
  freccia/.style   = {thick, -{Stealth}}
}

\title{Teorema di Eulero e potenze modulo $n$}
\author{Marco Cerbella}
\date{Lezione 14 -- 3 novembre 2026}

\begin{document}
\maketitle

\section{Calcolo di $\varphi(n)$}

\begin{teorema}[Corollario]
Sia $n = p_1^{h_1}\cdot p_2^{h_2}\cdots p_r^{h_r}$ la fattorizzazione in primi. Allora
\[
  \varphi(n) = \varphi(p_1^{h_1})\cdot\varphi(p_2^{h_2})\cdots\varphi(p_r^{h_r}) = \left(p_1^{h_1} - p_1^{h_1 - 1}\right)\left(p_2^{h_2} - p_2^{h_2 - 1}\right)\cdots\left(p_r^{h_r} - p_r^{h_r - 1}\right).
\]
\end{teorema}

\begin{esempio}
\begin{align*}
  \varphi(1000) &= \varphi(10^3) = \varphi(5^3\cdot 2^3) = \varphi(5^3)\,\varphi(2^3)\\
  &= (5^3 - 5^2)(2^3 - 2^2) = (125 - 25)(8 - 4) = 100\cdot 4 = 400,\\
  \varphi(360) &= \varphi(2^3\cdot 3^2\cdot 5) = \varphi(2^3)\,\varphi(3^2)\,\varphi(5)\\
  &= (2^3 - 2^2)(3^2 - 3)(5 - 1) = 4\cdot 6\cdot 4 = 96 .
\end{align*}
\end{esempio}

\section{Teorema di Eulero}

\begin{teorema}[di Eulero]
Sia $n > 1$ e $a \in \mathbb{Z}$. Se $\operatorname{MCD}(a, n) = 1$ allora
\[
  a^{\varphi(n)} \equiv 1 \pmod n .
\]
\end{teorema}

(Per $n = p$ primo, $\varphi(p) = p - 1$ e si ritrova il corollario del piccolo teorema di Fermat.)

\section{Come calcolare le potenze modulo $n$}

Si vuole calcolare $a^b \bmod n$.
\begin{enumerate}
  \item Se $a > n$, si riduce la base modulo $n$: $a = n\cdot q + a'$, e $a^b \equiv (a')^b \pmod n$.
  \item Se $\operatorname{MCD}(a', n) = 1$, si applica il teorema di Eulero: si calcola $\varphi(n)$ e si divide $b$ per $\varphi(n)$: $b = q\cdot\varphi(n) + r$. Allora
  \[
    a^b \equiv (a')^b \equiv (a')^{q\varphi(n) + r} = \left((a')^{\varphi(n)}\right)^q\cdot(a')^r \equiv (a')^r \pmod n,
  \]
  perché $(a')^{\varphi(n)} \equiv 1 \pmod n$.
\end{enumerate}

\begin{esempio}[$187^{5\,867\,432} \bmod 12$]
Si riduce la base modulo $12$: $187 = 12\cdot 15 + 7$, quindi $187^{5\,867\,432} \equiv 7^{5\,867\,432} \pmod{12}$. Si verifica che $\operatorname{MCD}(7, 12) = 1$, quindi si può applicare il teorema di Eulero:
\[
  \varphi(12) = \varphi(3\cdot 2^2) = \varphi(3)\,\varphi(2^2) = (3 - 1)(2^2 - 2) = 4 .
\]
Poiché $5\,867\,432$ è divisibile per $4$ (il resto è $0$), 
\[
  187^{5\,867\,432} \equiv 7^{5\,867\,432} \equiv 7^0 \equiv 1 \pmod{12} .
\]
\end{esempio}

\begin{esempio}[$4^{48\,642\,369} - 33^{1\,759\,468} \bmod 10$]
Qui $\operatorname{MCD}(4, 10) = 2$: i numeri non sono coprimi, quindi \emph{non} si può applicare il teorema di Eulero. Si cercano le potenze di $4$ modulo $10$:
\[
  \begin{gathered}
  4^0 = 1,\quad 4^1 = 4,\quad 4^2 = 16 \equiv 6,\\
  4^3 = 4^2\cdot 4 \equiv 6\cdot 4 = 24 \equiv 4,\quad 4^4 = 4^3\cdot 4 \equiv 4\cdot 4 = 16 \equiv 6 \pmod{10}.
  \end{gathered}
\]
Le potenze pari di $4$ sono $\equiv 6$, quelle dispari $\equiv 4 \pmod{10}$. L'esponente $48\,642\,369$ è dispari, quindi $4^{48\,642\,369} \equiv 4 \pmod{10}$.

Per l'altro termine: $33 \equiv 3 \pmod{10}$ e $\operatorname{MCD}(3, 10) = 1$, quindi si può applicare Eulero con $\varphi(10) = \varphi(2)\,\varphi(5) = (2 - 1)(5 - 1) = 4$. Poiché $1\,759\,468$ è divisibile per $4$ (resto $0$),
\[
  33^{1\,759\,468} \equiv 3^{1\,759\,468} \equiv 3^0 \equiv 1 \pmod{10}.
\]
In conclusione $4^{48\,642\,369} - 33^{1\,759\,468} \equiv 4 - 1 \equiv 3 \pmod{10}$.
\end{esempio}

\end{document}
